% from ega0-23-fr.tex
\oldpage[0\textsubscript{IV-1}]{221}

il suffira donc de prouver que $A'_{\mathfrak m}$ est l'intersection des
$A'_{\mathfrak p}$ pour les idéaux premiers $\mathfrak p\subset\mathfrak m$
de hauteur $1$. Comme $A'_{\mathfrak m}$ est la clôture intégrale de
$A_{\mathfrak m}$, on voit qu'on peut se borner à démontrer le corollaire pour
$A_{\mathfrak m}$. Or, il y a alors, en vertu de (23.2.7), une
$A_{\mathfrak m}$-algèbre finie $B$ contenue dans $A'_{\mathfrak m}$, telle
que $A'_{\mathfrak m}$ soit l'intersection des clôtures intégrales des anneaux
$B_{\mathfrak q}$ où $\mathfrak q$ parcourt l'ensemble des idéaux premiers de
hauteur $1$ de $B$. Mais $B$ est de la forme $C_{\mathfrak m}$, où $C$ est
une $A$-algèbre finie~; donc, pour tout idéal premier $\mathfrak q$ de hauteur
$1$ dans $B$, $\mathfrak q\cap A$ est par hypothèse un idéal premier de
hauteur $1$ dans $A$~; comme la clôture intégrale de $B_{\mathfrak q}$
contient celle de $A_{\mathfrak q\cap A}$ et que cette dernière contient
$A'_{\mathfrak m}$, il en résulte bien que $A'_{\mathfrak m}$ est
l'intersection des $A'_{\mathfrak p}$ pour $\mathfrak p\in P$ et
$\mathfrak p\subset\mathfrak m$, ce qui achève la démonstration.
\end{proof}

\begin{remarks}[23.2.10]
\label{0.23.2.10-fr}
\begin{enumerate}
  \item[(i)] Nous verrons (IV, 5.6.3 et 5.6.10) que l'hypothèse de (23.2.9)
  est vérifiée lorsque $A$ est \emph{universellement caténaire}, en particulier
  (IV, 5.10.17 et 5.11.2) lorsque $A$ est \emph{quotient d'un anneau
  régulier}~; enfin, elle est évidemment vérifiée si
  $\operatorname{Spec}(A')\to\operatorname{Spec}(A)$ est un homéomorphisme,
  en particulier si ce morphisme est radiciel (23.2.2).

  \item[(ii)] La conclusion de (23.2.9) n'est plus nécessairement exacte
  lorsqu'on omet l'hypothèse sur les sous-$A$-algèbres finies de $A'$. En
  effet, on peut définir un anneau local noethérien intègre $A$, de dimension
  $2$, dont la clôture intégrale $A'$ est une $A$-algèbre finie, mais tel qu'il
  y ait un idéal premier $\mathfrak p'$ de $A'$ de hauteur $1$ au-dessus de
  l'idéal maximal $\mathfrak m$ (de hauteur $2$) de $A$, et tel en outre que
  pour tout idéal premier $\mathfrak p$ de hauteur $1$ dans $A$,
  $A_{\mathfrak p}$ soit intégralement clos (IV, 5.6.11)~; l'intersection de
  ces anneaux $A_{\mathfrak p}$ ne peut alors être égale à $A'$ puisque $A'$
  est un anneau de Krull (Bourbaki, \emph{Alg. comm.}, chap.~VII, §~1,
  n\textsuperscript{o}~5, prop.~9).
\end{enumerate}
\end{remarks}

\begin{flushright}
\emph{(A suivre.)}
\end{flushright}
